\documentclass[11pt]{article}
% The preamble every project's paper shares: % The preamble every project's paper shares: % The preamble every project's paper shares: \input{../../preamble} from
% projects/<project>/paper/. Boards come from chessboard and xskak, plots from
% pgfplots; every number a paper prints comes from its generated.tex.
\usepackage[LSBC4,T1]{fontenc}
\usepackage[utf8]{inputenc}
\usepackage{lmodern}
\usepackage[margin=1in]{geometry}
\usepackage{microtype}
\usepackage{amsmath}
\usepackage{booktabs}
\usepackage{tabularx}
\usepackage{array}
\usepackage{xcolor}
\usepackage[obeyspaces,spaces]{url}
\usepackage{xskak}
\usepackage{pgfplots}
\pgfplotsset{compat=1.18}
\usepackage{caption}
\usepackage{subcaption}
\usepackage{enumitem}
\usepackage[section]{placeins}
\usepackage[hidelinks]{hyperref}

\setchessboard{boardfontencoding=LSBC4, boardfontfamily=skaknew, blackfieldcolor=black!22,
  whitefieldcolor=white, setfontcolors, showmover=false, label=true, labelfontsize=6pt,
  boardfontsize=14pt}

% A move or line in UCI, and a position in FEN, both breakable at a slash.
\newcommand{\uci}[1]{\texttt{#1}}
\newcommand{\fen}[1]{{\footnotesize\path{#1}}}
\newcolumntype{L}{>{\raggedright\arraybackslash}X}

\setlength{\parskip}{0.4em}
\setlength{\parindent}{0pt}
 from
% projects/<project>/paper/. Boards come from chessboard and xskak, plots from
% pgfplots; every number a paper prints comes from its generated.tex.
\usepackage[LSBC4,T1]{fontenc}
\usepackage[utf8]{inputenc}
\usepackage{lmodern}
\usepackage[margin=1in]{geometry}
\usepackage{microtype}
\usepackage{amsmath}
\usepackage{booktabs}
\usepackage{tabularx}
\usepackage{array}
\usepackage{xcolor}
\usepackage[obeyspaces,spaces]{url}
\usepackage{xskak}
\usepackage{pgfplots}
\pgfplotsset{compat=1.18}
\usepackage{caption}
\usepackage{subcaption}
\usepackage{enumitem}
\usepackage[section]{placeins}
\usepackage[hidelinks]{hyperref}

\setchessboard{boardfontencoding=LSBC4, boardfontfamily=skaknew, blackfieldcolor=black!22,
  whitefieldcolor=white, setfontcolors, showmover=false, label=true, labelfontsize=6pt,
  boardfontsize=14pt}

% A move or line in UCI, and a position in FEN, both breakable at a slash.
\newcommand{\uci}[1]{\texttt{#1}}
\newcommand{\fen}[1]{{\footnotesize\path{#1}}}
\newcolumntype{L}{>{\raggedright\arraybackslash}X}

\setlength{\parskip}{0.4em}
\setlength{\parindent}{0pt}
 from
% projects/<project>/paper/. Boards come from chessboard and xskak, plots from
% pgfplots; every number a paper prints comes from its generated.tex.
\usepackage[LSBC4,T1]{fontenc}
\usepackage[utf8]{inputenc}
\usepackage{lmodern}
\usepackage[margin=1in]{geometry}
\usepackage{microtype}
\usepackage{amsmath}
\usepackage{booktabs}
\usepackage{tabularx}
\usepackage{array}
\usepackage{xcolor}
\usepackage[obeyspaces,spaces]{url}
\usepackage{xskak}
\usepackage{pgfplots}
\pgfplotsset{compat=1.18}
\usepackage{caption}
\usepackage{subcaption}
\usepackage{enumitem}
\usepackage[section]{placeins}
\usepackage[hidelinks]{hyperref}

\setchessboard{boardfontencoding=LSBC4, boardfontfamily=skaknew, blackfieldcolor=black!22,
  whitefieldcolor=white, setfontcolors, showmover=false, label=true, labelfontsize=6pt,
  boardfontsize=14pt}

% A move or line in UCI, and a position in FEN, both breakable at a slash.
\newcommand{\uci}[1]{\texttt{#1}}
\newcommand{\fen}[1]{{\footnotesize\path{#1}}}
\newcolumntype{L}{>{\raggedright\arraybackslash}X}

\setlength{\parskip}{0.4em}
\setlength{\parindent}{0pt}

% Written by projects/review-walk/paper/generate.mjs; do not edit by hand.

\newcommand{\AppCommit}{e06307fe45}
\newcommand{\MaxErrors}{four}
\newcommand{\QuietErrors}{two}
\newcommand{\MaxHonors}{one}
\newcommand{\MateWithin}{three}
\newcommand{\LongWalk}{five}
\newcommand{\MaxStops}{five}
\newcommand{\HeldBefore}{2026-09-23}
\newcommand{\HeldAsked}{80}
\newcommand{\Draws}{10,000}
\newcommand{\MedianAbove}{50}
\newcommand{\QuartileBar}{40}
\newcommand{\TurnBar}{90}
\newcommand{\DesignGames}{734}
\newcommand{\DesignWithErrors}{707}
\newcommand{\DesignMomentsMedian}{6}
\newcommand{\DesignMomentsNinety}{11.7}
\newcommand{\DesignMomentsMost}{30}
\newcommand{\DesignStopsMedian}{3}
\newcommand{\DesignStopsNinety}{4}
\newcommand{\DesignStopsMost}{5}
\newcommand{\DesignLonger}{54}
\newcommand{\DesignNoStops}{19}
\newcommand{\DesignCoverage}{72}
\newcommand{\DesignCoverageLow}{70}
\newcommand{\DesignCoverageHigh}{73}
\newcommand{\DesignQuartile}{59}
\newcommand{\DesignTenth}{48}
\newcommand{\DesignSeriousGames}{636}
\newcommand{\DesignSeriousWhole}{71}
\newcommand{\DesignSeriousTenth}{77}
\newcommand{\DesignSeenIn}{707}
\newcommand{\DesignSeenOf}{707}
\newcommand{\DesignSeenAt}{63}
\newcommand{\DesignSeenAtLow}{28}
\newcommand{\DesignSeenAtHigh}{95}
\newcommand{\DesignTurnIn}{340}
\newcommand{\DesignTurnOf}{346}
\newcommand{\DesignTurnShare}{98}
\newcommand{\DesignMates}{27}
\newcommand{\DesignHonors}{173}
\newcommand{\DesignFirstStop}{39}
\newcommand{\DesignByCost}{87}
\newcommand{\DesignSoftened}{3}
\newcommand{\DesignFourthAdds}{4.6}
\newcommand{\DesignFifthAdds}{3.5}
\newcommand{\HeldGames}{80}
\newcommand{\HeldWithErrors}{79}
\newcommand{\HeldMomentsMedian}{6}
\newcommand{\HeldMomentsNinety}{11}
\newcommand{\HeldMomentsMost}{14}
\newcommand{\HeldStopsMedian}{3}
\newcommand{\HeldStopsNinety}{4}
\newcommand{\HeldStopsMost}{5}
\newcommand{\HeldLonger}{54}
\newcommand{\HeldNoStops}{1}
\newcommand{\HeldCoverage}{76}
\newcommand{\HeldCoverageLow}{70}
\newcommand{\HeldCoverageHigh}{79}
\newcommand{\HeldQuartile}{64}
\newcommand{\HeldTenth}{55}
\newcommand{\HeldSeriousGames}{70}
\newcommand{\HeldSeriousWhole}{71}
\newcommand{\HeldSeriousTenth}{83}
\newcommand{\HeldSeenIn}{79}
\newcommand{\HeldSeenOf}{79}
\newcommand{\HeldSeenAt}{59}
\newcommand{\HeldSeenAtLow}{30}
\newcommand{\HeldSeenAtHigh}{93}
\newcommand{\HeldTurnIn}{37}
\newcommand{\HeldTurnOf}{37}
\newcommand{\HeldTurnShare}{100}
\newcommand{\HeldMates}{4}
\newcommand{\HeldHonors}{14}
\newcommand{\HeldFirstStop}{38}
\newcommand{\HeldByCost}{88}
\newcommand{\HeldSoftened}{1}
\newcommand{\HeldFourthAdds}{6.7}
\newcommand{\HeldFifthAdds}{3.9}
\newcommand{\TeachStops}{212}
\newcommand{\TeachCards}{212}
\newcommand{\TeachCardShare}{100}
\newcommand{\TeachReply}{115}
\newcommand{\TeachReplyShare}{54}
\newcommand{\TeachChance}{50}
\newcommand{\TeachChanceShare}{24}
\newcommand{\TeachNeither}{47}
\newcommand{\TeachNeitherShare}{22}
\newcommand{\TeachLesson}{199}
\newcommand{\TeachLessonShare}{94}
\newcommand{\Kept}{kept}
\newcommand{\CheckMedian}{holds}
\newcommand{\CheckQuartile}{holds}
\newcommand{\CheckTurn}{holds}
\newcommand{\CheckLength}{holds}
\newcommand{\CheckSeen}{holds}
\newcommand{\MeasureTable}{%
\begin{table}[h]\centering\small
\begin{tabular}{lrr}
\toprule
 & design set & held out \\
\midrule
games & 734 & 80 \\
his marked moves a game: median, 90th, most & 6, 11.7, 30 & 6, 11, 14 \\
the walk's stops: median, 90th, most & 3, 4, 5 & 3, 4, 5 \\
games with more than five marked moves & 54\% & 54\% \\
cost held, median game (95\% interval) & 72\% (70 to 73) & 76\% (70 to 79) \\
cost held, lower-quartile game & 59\% & 64\% \\
a stop inside the turn, of games that turned on his moves & 340 of 346 & 37 of 37 \\
the costliest error is a stop & 707 of 707 & 79 of 79 \\
mistakes, misses and blunders all held & 71\% of 636 & 71\% of 70 \\
\bottomrule
\end{tabular}
\caption{The walk on the design set (the archive's reviews) and on the held-out phone reviews. Cost held is the \texttt{lossPct} of the errors the walk stops at over that of all his errors, inaccuracies included.}
\label{tab:measures}
\end{table}
}
\newcommand{\CapFigure}{%
\begin{figure}[h]\centering
\begin{tikzpicture}
\begin{axis}[width=0.78\linewidth, height=5.2cm, xmin=0.7, xmax=6.3, ymin=0, ymax=100, xtick={1,2,3,4,5,6}, xlabel={errors the walk may stop at}, ylabel={cost held (\%)}, legend pos=south east, legend cell align=left, legend style={font=\small, draw=none}, ymajorgrids, grid style={black!12}]
\addplot[thick, black, mark=*, mark size=1.6pt] coordinates {(1,37.1) (2,57.4) (3,67.0) (4,71.6) (5,75.1) (6,76.6)};
\addlegendentry{design set, median game}
\addplot[thick, black, dashed, mark=o, mark size=1.6pt] coordinates {(1,26.6) (2,43.5) (3,55.1) (4,59.2) (5,62.0) (6,63.9)};
\addlegendentry{design set, lower quartile}
\addplot[thick, red!70!black, mark=square*, mark size=1.5pt] coordinates {(1,35.4) (2,56.7) (3,68.9) (4,75.6) (5,79.5) (6,80.7)};
\addlegendentry{held out, median game}
\addplot[thick, red!70!black, dashed, mark=square, mark size=1.5pt] coordinates {(1,27.2) (2,46.5) (3,57.2) (4,64.3) (5,67.4) (6,68.7)};
\addlegendentry{held out, lower quartile}
\draw[black!45, densely dotted] (axis cs:4,0) -- (axis cs:4,100);
\end{axis}
\end{tikzpicture}
\caption{What the walk would hold of what his errors cost if it stopped at one to six of them, ranked as the walk ranks them. The dotted line is the cap the app uses.}
\label{fig:caps}
\end{figure}
}


\title{How many of his marked moves should a Game Review walk through?\\[0.3em]
\large Four costliest errors and one best move, from an overview}
\author{Chess Coach research lab}
\date{}

\begin{document}
\maketitle

\begin{abstract}
The app's Game Review marked every move of the player's that it graded an inaccuracy, a mistake,
a miss or a blunder, every great or brilliant move and every mate he let slip, and its walk
stopped at each: \DesignMomentsMedian{} stops in the median archived game, \DesignMomentsNinety{}
at the 90th percentile and \DesignMomentsMost{} at most. It also opened on his costliest error,
which falls anywhere from \DesignSeenAtLow\% to \DesignSeenAtHigh\% of the way into a game. The
change opens a review on an overview of the whole game and walks at most \MaxStops{} stops: his
\MaxErrors{} costliest mistakes, misses and blunders and \MaxHonors{} best move. On
\HeldGames{} held-out games reviewed as the phone reviews them, the walk's errors hold
\HeldCoverage\% of what all his errors cost in the median game (95\% interval \HeldCoverageLow\%
to \HeldCoverageHigh\%) and \HeldQuartile\% in the lower-quartile game, and a stop lies inside
the turn in \HeldTurnIn{} of the \HeldTurnOf{} games that turned on his own moves. By the rule
fixed before the run, the walk is \Kept.
\end{abstract}

\section{The question}

A review that is walked to its end teaches more than one that is abandoned, and the old walk was
long: more than \LongWalk{} stops in \DesignLonger\% of the archived games. It also had no fixed
place to begin. The screen opened on the move the player should have seen, his costliest error,
so each game opened on a different position, with the story of the game one tap away. The
question is what a short walk gives up. If the few costliest errors hold most of what his errors
cost, and one of them lies where the game turned, the walk can be short and still be the review
of the game.

\section{What the app does}

\paragraph{The walk.} The review service's \texttt{walkOf} reads a review's marked moves and
picks the stops. His mistakes, misses and blunders are ranked by what they cost (the review's
\texttt{lossPct}, the fall in his winning chances), the earlier first on a tie; then come the
mates within \MateWithin{} moves that he let slip, the shortest first; \MaxErrors{} of these in
all. When he made no mistake, miss or blunder, his \QuietErrors{} costliest inaccuracies stand
in. Then \MaxHonors{} best move of his joins them, a brilliant move before a great one. An inaccuracy never
takes a mistake's place, whatever it cost: the review grades an error an inaccuracy when the
position stayed about as good, and that grade is the review's own judgment of how much the move
mattered. The stops are walked in game order, and the other marked moves keep their marks in the
move line and on the chances graph.

\paragraph{The overview.} A review opens on the game as a whole: where it turned
(\texttt{turnOf}), the stops as a list, and what to take into the next game, written from the
move he should have seen (\texttt{seenOf}, which now ranks as the walk does, so it is always the
walk's costliest error). The chances graph numbers the stops. One button starts the walk at the
first stop he has not looked at.

\paragraph{What a stop says.} Each error's lesson card names a habit when the board code finds
one, and the list says which side of the board the habit is about: he missed their reply, or he
missed his own chance. Under the engine's lines, a stop links the lesson chapter that shows the
same mistake, when the study has one.

\section{The games and the rule}

The design set is the app's archive of his games as the routine reviewed them, \DesignGames{}
games, every one read while the walk was designed, with phone reviews of 37 recent games read
one by one. Held out are \HeldAsked{} games drawn with a fixed seed from the games played before
\HeldBefore, each reviewed anew the way the phone's Game Review reviews it, at the app commit
\texttt{\AppCommit}. No phone review of those games existed while the walk was designed. Their
archive reviews are in the design set's numbers: what is held out is the phone's review of them,
a smaller analysis with its own grades, and the one the screen walks.

The rule was committed before any held-out game was reviewed. The walk is kept when, on the
held-out games, its errors hold more than \MedianAbove\% of what his errors cost in the median
game, with the whole 95\% interval of that median above \MedianAbove\% (\Draws{} bootstrap
resamples of games), and at least \QuartileBar\% in the lower-quartile game; a stop lies inside
the turn in at least \TurnBar\% of the games that turned on his own moves; and two guards hold
that the rule's caps should make true of any game: no walk longer than \LongWalk{} stops with a
median of at most \MaxErrors, and the costliest error a stop in every game that has an error.

\section{Results}

\MeasureTable

Table~\ref{tab:measures} has both sets. Held out, the median bar \CheckMedian, the
lower-quartile bar \CheckQuartile{} and the turn bar \CheckTurn; the length guard \CheckLength{}
and the costliest-error guard \CheckSeen. The walk is \Kept.

In the median held-out game the phone marks \HeldMomentsMedian{} moves, as the routine does in
the median archived game (\DesignMomentsMedian), and the walk stops at \HeldStopsMedian. What the
walk removes most is the long tail: a tenth of the held-out reviews mark \HeldMomentsNinety{}
moves or more, and the longest \HeldMomentsMost, against a walk of at most \HeldStopsMost.

\CapFigure

\paragraph{The cap.} Four errors was a design choice, for a walk short enough to finish on a
phone, not a fitted value. Figure~\ref{fig:caps} shows what it gives up. On the design set the
fourth error adds \DesignFourthAdds{} points to the median game and a fifth would add
\DesignFifthAdds; held out, \HeldFourthAdds{} and \HeldFifthAdds. The curve flattens well below
the whole because the rest of the cost is in inaccuracies, which the walk leaves to the move
line whenever the game has a mistake: of his mistakes, misses and blunders alone, the walk holds
every one in \HeldSeriousWhole\% of the \HeldSeriousGames{} held-out games that have one
(\DesignSeriousWhole\% of \DesignSeriousGames{} in the design set). Four stops ranked by cost
alone, whatever the grade, and always filled to four, would hold \HeldByCost\% in the median
held-out game, and would put a softened error ahead of a mistake in \HeldSoftened\% of the
games.

\paragraph{Where the old screen opened.} The costliest error comes \HeldSeenAt\% of the way into
the median held-out game, and anywhere from \HeldSeenAtLow\% to \HeldSeenAtHigh\% between the
lowest and highest tenth (\DesignSeenAtLow\% to \DesignSeenAtHigh\% in the design set). A screen
that opens on it opens on a different part of every game.

\paragraph{What the stops say.} Of the walk's \TeachStops{} held-out error stops, \TeachCards{}
(\TeachCardShare\%) have a lesson card. \TeachReply{} of the cards (\TeachReplyShare\%) say he
missed their reply, \TeachChance{} (\TeachChanceShare\%) that he missed his own chance, and
\TeachNeither{} (\TeachNeitherShare\%) name neither. \TeachLesson{} of the stops
(\TeachLessonShare\%) link a lesson chapter. These are counts of what the screen says, not a judgment of whether it is the right thing to say.

\section{Limits}

Cost is the engine's reading of winning chances, so the walk stops at what cost most by the
engine, not at what he could most easily have found. Ranking the stops by how findable the
better move was at his level is the next question, and it needs the human-move model on every
stop. The held-out games are one player's, and the same games' archive reviews were part of the
design set. The two guards are true by construction and are checked so that a later change to
the caps cannot pass unseen.

\section{What the app took}

The walk, the overview a review opens on, and \texttt{seenOf}'s ranking are in the chess pull
request this project measured, at \texttt{\AppCommit}. The pull request's later commits change
the screen's layout and wording, not which moves the walk stops at or which chapter a stop
links.

\end{document}
